\section{One actual polydisc for all canonical normal coefficient functions}

\begin{theorem}[Original canonical generators give all coefficient functions on one true domain]
\label{hir:actual-uniform-canonical-coefficient-functions}
\uses{hir:actual-canonical-reduced-analytic-coefficient-data,hir:actual-original-finite-family-coefficient-radius,hir:actual-relative-coefficient-common-polydisc-functions,hir:actual-convergent-original-taylor-realization}
\lean{Hironaka.HironakaActualUniformCanonicalCoefficients.exists_uniform_canonical_coefficient_data}
\leanok
For every original whole germ ideal $I$ in the actual coordinate blocks
$(u,z)$, internally construct one actual analytic vertex coefficient datum
$P$. Its original generators have exactly the canonical reduced Taylor
series of the same whole formal extension and generate the same original
ideal. Every coefficient function of this same datum is holomorphic on
its single actual open neighbourhood of zero, before all normal labels.
Each coefficient function of this chosen datum is the literal convergent evaluation of the
corresponding relative coefficient of that same canonical reduced series.
\end{theorem}
\begin{proof}
Choose the genuine original canonical analytic reduced family once.
The finite-family original Taylor mass theorem, transported through the
literal dimension equality $0+n=n$, supplies one positive uniform full
radius for all unchanged generator series. The genuine coefficient mass
slicing theorem supplies the same base radius for every relative normal
coefficient of every generator. Evaluate these actual relative series
in the original base coordinates. All of these functions are holomorphic
on the same actual open base polydisc, which contains zero.

The genuine inverse Taylor realization theorem, transported through that
same dimension equality, proves that the original Taylor germ of each
such actual evaluation is exactly its unchanged relative coefficient
series. Thus the original coefficient Taylor identities hold for these
chosen functions. Keep the actual original generators, whole-ideal
generation, initial vertices and monicity. Use the same true base
polydisc and these actual functions as the coefficient fields of $P$.
Their actual holomorphy proves all required low-coefficient regularity,
coherence of their genuine principal ideals, and membership of the
actual function in its own ideal. This constructs the complete genuine
datum and the stronger all-label common-domain assertion internally.
\end{proof}

\begin{theorem}[All whole-diagram coefficient identities hold as functions on one common domain]
\label{hir:actual-uniform-canonical-diagram-function-identities}
\uses{hir:actual-uniform-canonical-coefficient-functions,hir:actual-initial-tangent-stable-pure-normal-vertices,hir:actual-pure-normal-reduced-diagram-coefficients}
\lean{Hironaka.HironakaActualUniformCanonicalCoefficients.exists_uniform_canonical_diagram_data}
\leanok
Suppose the same whole homogeneous initial ideal is stable under the
actual tangent partial derivatives. One internally produced datum $P$
retains the preceding exact canonical generator Taylor identities and
all-label common-domain holomorphy and, simultaneously for every normal
label $(0,\gamma)$ in its actual whole diagram and every actual point
of its same domain,
\[
 P_{a,\gamma}(u)=
 \begin{cases}1&\gamma=a_z,\\0&\gamma\ne a_z.\end{cases}
\]
No finite cutoff or label-dependent shrinking is needed.
\end{theorem}
\begin{proof}
Construct the one actual datum by the preceding theorem. Genuine tangent
initial stability makes every actual vertex pure normal. Upperness of
the whole diagram and the genuine reduced coefficient formula identify
the entire formal relative coefficient at every diagram label with its
literal constant zero or one series. The actual coefficient function
chosen above is exactly the convergent evaluation of that entire series.
Evaluating the constant series gives the same constant at every point
of the same domain. The radius and datum were chosen before every label,
so the resulting function identities hold simultaneously. No infinitely
many independent neighbourhood identities have been intersected.
\end{proof}

\begin{theorem}[Original constant-HS centre geometry gives uniform canonical analytic functions]
\label{hir:actual-original-hs-centre-uniform-canonical-functions}
\uses{hir:actual-uniform-canonical-diagram-function-identities,hir:actual-original-coupled-normal-free-initial-cylinder,hir:actual-original-centre-reduced-normal-powers}
\lean{Hironaka.HironakaActualUniformCanonicalCoefficients.exists_original_hs_centre_uniform_canonical_data}
\leanok
An original regular centre, original coherent ideal, and all-order
original Hilbert--Samuel constancy produce one actual flattening and
one actual canonical coefficient datum there. All coefficient functions
are holomorphic on one actual common base neighbourhood and all preceding
whole-diagram identities hold there simultaneously as function identities.
The same original canonical generators lie in their actual prescribed
whole normal powers. The original chart source, origin and centre
identities are retained. No chart, family, radius, coefficient-support
or normal-power certificate is supplied.
\end{theorem}
\begin{proof}
Invoke the genuine coupled original normal-data producer once. It fixes
one actual original flattening and the same literal flattened origin
ideal, with whole tangent initial stability and actual normal-layer
freeness. Apply the preceding construction to that same ideal and
stability, obtaining one datum and all uniform function identities.
The coupled stability and freeness give the actual canonical reduced
normal powers. Each chosen unchanged original generator has exactly
that reduced Taylor series; the genuine analytic/formal normal-power
dictionary therefore puts it in the same actual whole normal power.
Retain the original actual chart identities and this one produced datum.
\end{proof}

This constructs genuine uniform coefficient functions. A current-history
Hilbert--Samuel persistence equivalence, normalized-versus-geometric
strict-ideal equality and a terminating driver remain separate obligations.
